\newcommand{\figuraCurvaturasPrincipales}{%
\begin{figure}[H]
\centering
\begin{tikzpicture}[>=Latex,font=\small,line cap=round,line join=round]
  \node[font=\bfseries] at (3,5.0) {(a) Curva plana};
  \draw[very thick,blue!70!black] (0.6,1.1)
    .. controls (2.0,4.7) and (4.2,4.5) .. (5.4,1.3);
  \coordinate (P) at (3.0,3.75);
  \fill (P) circle (0.06) node[above right] {$P$};
  \draw[thick,dashed] (1.4,3.2)--(4.6,4.3) node[right] {tangente};
  \draw[thick,densely dashed,red!70!black] (3.0,2.05) circle (1.70);
  \draw[->,red!70!black] (3.0,2.05)--(P) node[midway,right] {$R$};

  \node[font=\bfseries] at (9,5.0) {(b) Superficie};
  \fill[blue!10] (6.4,1.1)
    .. controls (7.0,4.5) and (10.7,4.5) .. (11.6,1.2)
    .. controls (10.2,0.3) and (7.7,0.3) .. cycle;
  \draw[very thick,blue!70!black] (6.4,1.1)
    .. controls (7.0,4.5) and (10.7,4.5) .. (11.6,1.2);
  \coordinate (Q) at (9.0,3.4);
  \fill (Q) circle (0.06) node[above] {$P$};
  \draw[->,thick] (Q)--++(90:1.25) node[above] {$\hat n$};
  \draw[<->,thick,dashed,red!70!black] ($(Q)+(-1.7,0)$)--($(Q)+(1.7,0)$)
    node[right] {dirección principal 1};
  \draw[<->,thick,dashed,green!45!black] ($(Q)+(0,0.95)$)--($(Q)+(0,-0.95)$)
    node[below] {dirección principal 2};
  \draw[very thick] ($(Q)+(-0.35,-0.22)$)--++(0.75,0.12)--++(0.25,0.5)--++(-0.75,-0.12)--cycle;
\end{tikzpicture}
\caption{Geometría local de la curvatura. Una curva se aproxima mediante su
círculo osculador; en una superficie, las curvaturas extrema y mínima definen
las dos direcciones principales.}
\label{fig:curvaturas-principales}
\end{figure}%
}